\begin{lemma}
Under the hypotheses of the opposite-projection collision bound, the probability that the size bound holds but the red estimate fails is $o(1)$. Hence the conditional red estimate holds with high probability.
\end{lemma}
\begin{proof}
Write \(M=\noderef{TwoBiteNaturalReducedVertexCount}(n)\),
\(p=\noderef{TwoBiteEdgeProbability}(1/2,n)\), and \(L=\log n\).  Fix a
red ordered pair \(r\ne s\).  Put
\[
U=N_r,\qquad V=N_s,\qquad U_0=U\setminus V,\qquad V_0=V\setminus U.
\]
Condition on the red base graph and the red row map \(\rho=\pi_R\).  The
conditioning cells form a finite partition, and
\(\noderef{OppositeProjectionConditioning}\) allows us to prove a uniform
bound on each nonzero cell; cells of zero mass contribute nothing because
\(\noderef{TwoBiteConditionalProbability}\) is defined to be \(0\) there.

On the size event,
\[
|U_0|\le |U|\le 2pn,\qquad |V_0|\le |V|\le 2pn.
\]
After the red row map is fixed, the remaining blue coordinates are uniform
row injections.  This is exactly the row law in
\(\noderef{OppositeProjectionRowInjectionLaw}\): if \(T\subseteq V_0\) is
fixed and \(S\subseteq V_B\), then the conditional probability that every
vertex of \(T\) lands in \(S\) is at most \((|S|/M)^{|T|}\), with the
zero-mass case already included.

Now expose all blue coordinates of \(U_0\).  Let \(\eta\) denote a possible
full exposure of these coordinates, and write
\[
S(\eta)=\pi_B(U_0),\qquad |S(\eta)|\le |U_0|.
\]
The rows used by \(U_0\) and \(V_0\) are disjoint: \(U_0\) lies over red
rows adjacent to \(r\) but not \(s\), while \(V_0\) lies over red rows
adjacent to \(s\) but not \(r\).  The sets are also disjoint by definition of
the set differences.  Thus the hypotheses of
\(\noderef{OppositeProjectionExposureTailBound}\) hold for every exposure
cell \(\eta\).

For a fixed exposure \(\eta\), define
\[
W_\eta=|\{v\in V_0:\pi_B(v)\in S(\eta)\}|.
\]
The exposure-tail helper first handles the zero-mass exposure cells using the
zero convention in \(\noderef{TwoBiteConditionalProbability}\).  In every
nonzero cell it expands the finite sample weights, cancels the row factors
fixed by the \(U_0\)-exposure, and uses the unchanged row-injection law on
the rows meeting \(V_0\).  Consequently, for every \(t\),
\[
\Pr(W_\eta\ge t\mid G_R,\rho,\eta)
\le \binom{|V_0|}{t}\left(\frac{|S(\eta)|}{M}\right)^t
\le \binom{|V_0|}{t}\left(\frac{|U_0|}{M}\right)^t .
\]
Average this inequality over all possible exposures \(\eta\).  The exposure
events form a finite partition of the already conditioned cell; each has
nonnegative mass, and zero-mass exposure events contribute \(0\).  Since the
right hand side is uniform in \(\eta\), the same bound holds before exposure:
\[
\Pr(W\ge t\mid G_R,\rho)
\le \binom{|V_0|}{t}\left(\frac{|U_0|}{M}\right)^t
\le \left(\frac{e|U_0||V_0|}{Mt}\right)^t.
\]
The last inequality is the standard estimate
\(\binom bt\le(eb/t)^t\), and the event is empty when \(t>|V_0|\).

The intersection-containment node
\(\noderef{OppositeProjectionIntersectionContainment}\) gives
\[
|\pi_B(U)\cap\pi_B(V)|
\le |U\cap V|+W.
\]
Indeed, the only projected intersection points not already coming from
\(U\cap V\) are witnessed by a vertex of \(V_0\) whose blue coordinate lies in
\(\pi_B(U_0)\), and these witnesses are counted by \(W\).

Now estimate the parameter in the tail bound.  Since
\[
p=\noderef{TwoBiteEdgeProbability}(1/2,n)
=\frac12\sqrt{\frac{L}{n}},
\]
and \(L\ge0\) for all sufficiently large \(n\), we have
\[
p^2n=\frac14\frac{L}{n}\,n=\frac{L}{4}.
\]
Also
\[
M=\left\lfloor\frac{n}{L^2}\right\rfloor .
\]
Because \(n/L^2\to\infty\), for all sufficiently large \(n\) we have
\(n/L^2\ge2\), and then
\[
M=\left\lfloor\frac{n}{L^2}\right\rfloor
\ge \frac12\frac{n}{L^2}.
\]
Indeed, for \(x\ge2\), \(\lfloor x\rfloor\ge x-1\ge x/2\).  Combining this
with the size bounds gives, for all sufficiently large \(n\),
\[
\frac{|U_0||V_0|}{M}
\le \frac{(2pn)^2}{M}
=\frac{4p^2n^2}{M}
\le \frac{Ln}{M}
\le 2L^3.
\]
Take \(t=\lceil100L^3\rceil\).  Then \(t\ge100L^3\), so
\[
\frac{e|U_0||V_0|}{Mt}
\le \frac{2eL^3}{100L^3}
=\frac e{50}.
\]
Since \(e/50<1\), the tail estimate gives
\[
\left(\frac{e|U_0||V_0|}{Mt}\right)^t
\le \left(\frac e{50}\right)^t
\le \left(\frac e{50}\right)^{100L^3}
=\exp(-cL^3),
\]
where \(c=100\log(50/e)>0\).  If \(t>|V_0|\), the tail event is empty, so the
same upper bound holds trivially.  Thus a fixed red pair fails with
probability at most \(\exp(-cL^3)\), uniformly over the conditioning data.

There are at most \(M^2\le n^2=\exp(2L)\) ordered red pairs for all large
\(n\).  The union bound gives
\[
M^2\exp(-cL^3)\to0
\]
by \(\noderef{OppositeProjectionAsymptoticBound}\).  Therefore, with high
probability, no red pair violates the desired projected intersection estimate
while the size bound holds.
\end{proof}
